% Nouveau EGA V (cote 25): archive cote 25, site page 32 (PDF page 33).
% typescript. Draft compared: none.
% AI-transcribed by Claude Opus 5.5 (claude-opus-5-5, adaptive thinking, effort high), Message Batches API, batch ega25b (msgbatch_01A7n5YuXubUPNK1zsepi8RG), 2026-11-13; original AI transcription, subsequently checked against both supplied scan parts; see PAGES_CHECKED.tsv and CHECK_DECISIONS.tsv.
\scanpage{25/32}
\mspage{\hand{28}}

8.2. Soit Z l'ensemble des $\hand{\xi\in}\overset{\hand{\vee}}{P}$ tels que dim $Y_{\hand{\xi}}$ $\hand{>}$ dim X - 1,\hand{\unsure{\xi}}
et supposons que pour toute composante irr $X_i$ de X, on ait dim $f(X_i)$ $\hand{>}$ 0;
alors Z est de codimension 2 dans $\overset{\hand{\vee}}{P}$.  Cela
2.1. et 2.2.
résulte de 
 (qui implique que toute comp irr de $\underline{Y}$ domine $\overset{\hand{\vee}}{P}$ ) et de
la théorie de la dimension pour le morphisme $\underline{Y} \hand{\to} \overset{\hand{\vee}}{P}$ . A partir de ce
résultat  on peut donner comme corollaire le cas où on
part d'une partie fermée  Z de X et qu'on r egarde la dimension des
images inverses $Z_{\hand{\xi}}$ dans les $Y_{\hand{\xi}}$ ( $\hand{\xi\in}\overset{\hand{\vee}}{P}$ ), et on peut même prendre pour
Z l'ensemble des $\hand{\xi\in}\overset{\hand{\vee}}{P}$ tels que il existe une co mposante irr. de $T_{k(\hand{\xi})}$
dont la trace sur $Y_{\hand{\xi}}$ soit de dimension trop grande (NB on suppose toujour
que pour toute comp irréd $Z_i$ de Z, on ait dim $f(Z_{\hand{i}})$ $\hand{>}$ 0). Enfin, l'énoncé
le plus précis dans cette direction, et qui résulte aisément du premier
énoncé (pour X irréductible) et de 2.7. , est l'énoncé modulé suivant:
F étant cohérent sur X,  supposons que pour tout cycle premier associ
T pour F, on ait dim f(T)$\hand{>}$ 0\unsure{;|,} alors l'ensemble des $\hand{\xi\in}\overset{\hand{\vee}}{P}$ tels que
$\hand{\phi_\xi}$ ne soit pas $F_{k(\hand{\xi})}$-régulier est (constructible et) de codimension$\hand{\geqslant}$ 2.
\placedmargin{\textbar} (La notation pour $\hand{\phi_\xi}$ est celle du n°5). On peut donner ceci comme
\placedmargin{\textbar} proposition principale, et énoncer les  assertions qui précédaient
\placedmargin{\textbar} comme corollaires, la démonstration procédant via l'un des corollaires.



\mspage{\hand{32}}
